\documentclass[11pt]{article}
\usepackage[T1]{fontenc}
\usepackage[utf8]{inputenc}
\usepackage{amsmath,amssymb,amsthm}
\usepackage[margin=1in]{geometry}
\usepackage[colorlinks=true,linkcolor=blue,citecolor=blue,urlcolor=blue]{hyperref}
\newtheorem{theorem}{Theorem}
\newtheorem{lemma}{Lemma}
\newcommand{\Q}{\mathbb{Q}}
\newcommand{\R}{\mathbb{R}}
\newcommand{\wt}{\operatorname{wgt}_{\mathrm{TC}}}
\newcommand{\genus}{\operatorname{genus}}
\title{A positive-norm counterexample to a TC-weight question}
\author{Kaiwen Jin\\\small\href{mailto:yui.ui.siki@gmail.com}{yui.ui.siki@gmail.com}}
\date{}

\begin{document}
\maketitle

\begin{abstract}
Let $S$ be the closed oriented surface of genus two. The aspherical
four-manifold $M=S\times S$ has a rational fundamental class of positive
homological $\ell^1$ seminorm, but every nonzero $u\in H^4(M;\Q)$ satisfies
$\wt(1\times u-u\times1)\le2$. This gives a negative answer to
Neofytidis's question asking whether a positive-norm degree-$k$ class on a
closed oriented aspherical manifold forces the existence of a nonzero
degree-$k$ class with TC weight at least $k$.
\end{abstract}

We use the unreduced convention for genus: $\genus(p)$ is the least number
of open sets covering the base of $p$ on which continuous local sections
exist. Let
\[
 e_M:C([0,1],M)\longrightarrow M\times M,
 \qquad e_M(\gamma)=(\gamma(0),\gamma(1)),
\]
where the path space has the compact-open topology. For a nonzero class
$v\in H^*(M\times M;\Q)$ its TC weight is
\begin{equation}\label{eq:weight}
 \wt(v)=\sup\bigl\{r\in\mathbb Z_{\ge0}:
 F^*v=0\ \text{for every }F:Y\to M\times M
 \text{ with }\genus(F^*e_M)\le r\bigr\}.
\end{equation}
Here $Y$ ranges over all topological spaces and $F$ is continuous. This is
the sectional-weight convention of Farber--Grant
\cite[Definition 29]{FG}, as used in \cite[Definition 2.4]{N}.
For $u\in H^k(M;\Q)$ write
$\bar u=\operatorname{pr}_2^*u-\operatorname{pr}_1^*u
=1\times u-u\times1$.

For $\alpha\in H_k(M;\Q)$, the rational homological seminorm is
\[
 \|\alpha\|_{1,\Q}
 =\inf\left\{\sum_{j=1}^m|a_j|:
 \sum_{j=1}^m a_j\sigma_j\text{ is a rational singular cycle
 representing }\alpha\right\}.
\]
Neofytidis asks whether, for every closed connected oriented aspherical
$n$-manifold $M$ and $2\le k\le n$, the existence of an
$\alpha\in H_k(M;\Q)$ with $\|\alpha\|_{1,\Q}>0$ implies the
existence of a nonzero $u\in H^k(M;\Q)$ with $\wt(\bar u)\ge k$
\cite[Question 5.1(a)]{N}. The following example refutes that implication.

\begin{theorem}\label{thm:main}
Let $S$ be the closed oriented surface of genus two and let $M=S\times S$,
with its product orientation. Then $M$ is a closed connected oriented
aspherical four-manifold and
\[
 \|[M]_\Q\|_{1,\Q}\ge24>0.
\]
For every nonzero $u\in H^4(M;\Q)$,
\[
 \wt(\bar u)\le2<4.
\]
In particular, the assertion in the question fails for $n=k=4$ and
$\alpha=[M]_\Q$.
\end{theorem}

The weight bound comes from keeping the first surface coordinate fixed
while moving the second to one base point. We first give the required
open cover and its local sections explicitly.

\begin{lemma}\label{lem:sections}
Fix $b\in S$, put $h(x,y)=(x,b)$, and define
$F:M\to M\times M$ by $F(z)=(z,h(z))$. Then
$\genus(F^*e_M)\le3$.
\end{lemma}

\begin{proof}
Choose a finite simplicial triangulation $K$ of the surface. In its
barycentric subdivision, a vertex $b_\tau$ corresponds to a nonempty
simplex $\tau$ of $K$. Color $b_\tau$ by $\dim\tau\in\{0,1,2\}$ and put
\[
 U_i=\bigcup_{\dim\tau=i}\operatorname{st}^{\circ}(b_\tau),
 \qquad i=0,1,2.
\]
The open star of a vertex is the set where its barycentric coordinate
is positive. The sets $U_i$ are open and cover $S$. Distinct vertices of
the same color do not belong to one simplex of the subdivision, because
each such simplex corresponds to a strictly increasing chain of
simplices of $K$. Their open stars are therefore disjoint.

Each open star contracts to its vertex by the straight-line homotopy
within simplices containing that vertex. Append a fixed path from that
vertex to $b$. Since the stars of a fixed color are disjoint open sets,
these homotopies together give a continuous map
\[
 H_i:U_i\times[0,1]\longrightarrow S,
 \qquad H_i(y,0)=y,\quad H_i(y,1)=b.
\]
The three sets $V_i=S\times U_i$ cover $M$. On $V_i$ the path
\[
 \gamma_i(x,y)(t)=(x,H_i(y,t))
\]
has endpoints $(x,y)$ and $(x,b)$. The adjoint
$V_i\times[0,1]\to M$ is continuous, and hence
$\gamma_i:V_i\to C([0,1],M)$ is continuous in the compact-open topology.
Indeed, for a compact $C\subset[0,1]$ and an open $O\subset M$, if
$\gamma_i(z_0)(C)\subset O$, continuity of the adjoint gives product
neighborhoods of $(z_0,t)$ mapped into $O$ for each $t\in C$. A finite
subcover of $C$ and the intersection of the corresponding neighborhoods
of $z_0$ show that $\gamma_i^{-1}[C,O]$ is open.
Thus $z\mapsto(z,\gamma_i(z))$ is a continuous section of $F^*e_M$ on
$V_i$, proving the bound.
\end{proof}

\begin{proof}[Proof of Theorem~\ref{thm:main}]
The surface $S$ has a contractible universal cover, and the product of
two such covers is a contractible universal cover of $M$; see
\cite[Examples 1B.2, 1B.5 and 1B.14]{H}. Thus $M$ is aspherical.
The other manifold properties follow from the corresponding properties
of the oriented surface.

Bucher-Karlsson's surface-product formula
\cite[Corollary 3]{B} gives the real simplicial volume
\[
 \|[M]_\R\|_{1,\R}=24(2-1)(2-1)=24.
\]
Every rational singular cycle representing $[M]_\Q$ represents
$[M]_\R$ after coefficient extension and has the same $\ell^1$ chain
norm. Each such norm is therefore at least $24$, so their infimum
satisfies $\|[M]_\Q\|_{1,\Q}\ge24$. No equality between the two
coefficient-field seminorms is needed.

Use the maps $h$ and $F$ of Lemma~\ref{lem:sections}. The map $h$ factors
as $M\xrightarrow{q}S\xrightarrow{j}M$, where $q(x,y)=x$ and
$j(x)=(x,b)$. Since $S$ is a two-dimensional CW complex,
$H^4(S;\Q)=0$ \cite[Theorem 3.5]{H}, and hence
$h^*u=q^*j^*u=0$ for every $u\in H^4(M;\Q)$. Consequently
\begin{equation}\label{eq:pullback}
 F^*\bar u=h^*u-\operatorname{id}_M^*u=-u.
\end{equation}
For every nonzero $u$, this same map $F$ has nonzero pullback and
$\genus(F^*e_M)\le3$. The vanishing property in
\eqref{eq:weight} therefore fails for every integer $r\ge3$, giving
$\wt(\bar u)\le2$. Equation~\eqref{eq:pullback} also shows that
$\bar u\ne0$. Nonzero degree-four classes exist, since Poincar\'e
duality gives $H^4(M;\Q)\cong H_0(M;\Q)\cong\Q$
\cite[Theorem 3.30]{H}. Thus no choice of a nonzero class can satisfy
the requested weight bound.
\end{proof}

The local-section argument is a concrete instance of existing genus
machinery. For maps $f,g$ with the same domain and codomain,
$D(f,g)\le d$ means that $d+1$ open sets cover their domain and the
restrictions of $f,g$ are homotopic on each set, where $d$ is a
nonnegative integer. Mac\'ias-Virg\'os and Mosquera-Lois prove
\cite[Theorem 2.8]{D} that, for the pair $(\operatorname{id}_M,h)$,
\[
 \genus(F^*e_M)=D(\operatorname{id}_M,h)+1.
\]
Write $c_b:S\to S$ for the constant map at $b$, so that
$h=\operatorname{id}_S\times c_b$. Their product inequality
\cite[Theorem 3.19]{D} applies because $M=S\times S$ is normal and gives
\[
 D(\operatorname{id}_M,h)
 \le D(\operatorname{id}_S,\operatorname{id}_S)
       +D(\operatorname{id}_S,c_b)\le0+2.
\]
The last bound follows from the three open surface sets in the proof of
Lemma~\ref{lem:sections}, on which the identity is homotopic to $c_b$.
Thus the normalized distance bound gives the unreduced genus bound
three used above. Here the
surface-product volume formula and the degree-four pullback
\eqref{eq:pullback} connect this genus bound to the positive-norm
question. The weight estimate concerns the difference $\bar u$. For comparison,
the unreduced topological complexity is $\mathrm{TC}(M)=\genus(e_M)$;
for this surface product it equals $9$ \cite[Example 5.7]{N}.

\paragraph{AI use.}
OpenAI Codex was used to develop the argument, check cited source passages,
prepare the text, and conduct automated manuscript reviews. These reviews
do not constitute independent mathematical peer review.

\begin{thebibliography}{9}
\bibitem{B}
M.~Bucher-Karlsson,
The simplicial volume of closed manifolds covered by
$\mathbb H^2\times\mathbb H^2$,
\emph{Journal of Topology} \textbf{1} (2008), 584--602.
\href{https://doi.org/10.1112/jtopol/jtn012}{doi:10.1112/jtopol/jtn012}.
\href{https://arxiv.org/abs/math/0703587v2}{arXiv:math/0703587v2}.

\bibitem{FG}
M.~Farber and M.~Grant,
Symmetric motion planning,
\emph{Contemporary Mathematics} \textbf{438} (2007), 85--104.
\href{https://arxiv.org/abs/math/0610857}{arXiv:math/0610857}.

\bibitem{H}
A.~Hatcher,
\emph{Algebraic Topology}, Cambridge University Press.
The cited electronic text is marked copyright 2002:
\href{https://webhomes.maths.ed.ac.uk/~v1ranick/papers/hatchat.pdf}{consulted electronic version}.

\bibitem{D}
E.~Mac\'ias-Virg\'os and D.~Mosquera-Lois,
Homotopic distance between maps,
\emph{Mathematical Proceedings of the Cambridge Philosophical Society}
\textbf{172} (2022), 73--93.
\href{https://arxiv.org/abs/1810.12591v3}{arXiv:1810.12591v3}.

\bibitem{N}
C.~Neofytidis,
Topological complexity, asphericity and connected sums,
\emph{Journal of Applied and Computational Topology} \textbf{9} (2025),
article 10.
\href{https://doi.org/10.1007/s41468-025-00205-z}{doi:10.1007/s41468-025-00205-z}.
The question and numbered locators cited here refer to
\href{https://arxiv.org/abs/2212.08962v2}{arXiv:2212.08962v2}.
\end{thebibliography}
\end{document}
